\section{2018年高考题}
\subsection{2018年全国I卷（理）}\label{gk:2018年全国I卷（理）}


\clearpage


\subsection{2018年全国I卷（文）}\label{gk:2018年全国I卷（文）}


\clearpage


\subsection{2018年全国II卷（理）}\label{gk:2018年全国II卷（理）}


\clearpage


\subsection{2018年全国II卷（文）}\label{gk:2018年全国II卷（文）}


\clearpage


\subsection{2018年全国III卷（理）}\label{gk:2018年全国III卷（理）}


\clearpage


\subsection{2018年全国III卷（文）}\label{gk:2018年全国III卷（文）}


\clearpage


\subsection{2018年北京卷（理）}\label{gk:2018年北京卷（理）}

\begin{description}
    \item[一、]
    \begin{enumerate}
        \item 1
        \item 2
        \item 3
        \item 4
        \item 5
        \item 6
        \item 7
        \item 8
    \end{enumerate}
    \item[二、]
    \begin{enumerate}
      \addtocounter{enumi}{+8}   
        \item 1
        \item 2
        \item 3
        \item 4
        \item 1
        \item 6
    \end{enumerate}
    \item[三、]   
    \begin{enumerate}
      \addtocounter{enumi}{+14}   
        \item 1
        \item 2
        \item 3
        \item  设函数 $f(x)=[ax^{2}-(4a+1)x+4a+3]e^{x}$。

\begin{enumerate}[label=(\arabic*)]
    \item 若曲线 $y=f(x)$ 在点 $(1,f(1))$ 处的切线与 $x$ 轴平行，求 $a$；
    \item 若 $f(x)$ 在 $x=2$ 处取得极小值，求 $a$ 的取值范围。
\end{enumerate}
        \item 已知抛物线 $C: y^{2}=2px$ 经过点 $P(1,2)$。过点 $Q(0,1)$ 的直线 $l$ 与抛物线 $C$ 有两个不同的交点 $A,B$，且直线 $PA$ 交 $y$ 轴于 $M$，直线 $PB$ 交 $y$ 轴于 $N$。

\begin{enumerate}[label=(\arabic*)]
    \item 求直线 $l$ 的斜率的取值范围；
    \item 设 $O$ 为原点，$\overrightarrow{QM}=\lambda \overrightarrow{QO}$，$\overrightarrow{QN}=\mu \overrightarrow{QO}$，求证：$\frac{1}{\lambda}+\frac{1}{\mu}$ 为定值。
\end{enumerate}
        \item 设 $n$ 为正整数，集合 $A=\left\{\alpha \mid \alpha=\left(t_{1},t_{2},\cdots ,t_{n}\right),t_{k}\in\left\{0,1\right\},k=1,2,\cdots ,n\right\}$。对于集合 $A$ 中的任意元素 $\alpha=\left(x_{1},x_{2},\cdots ,x_{n}\right)$ 和 $\beta=\left(y_{1},y_{2},\cdots ,y_{n}\right)$，记 $M(\alpha,\beta)=\frac{1}{2}[(x_{1}+y_{1}-\left|x_{1}-y_{1}\right|)+(x_{2}+y_{2}-\left|x_{2}-y_{2}\right|)+\cdots +(x_{n}+y_{n}-\left|x_{n}-y_{n}\right|)]$。

\begin{enumerate}[label=(\arabic*)]
    \item 当 $n=3$ 时，若 $\alpha=(1,1,0)$，$\beta=(0,1,1)$，求 $M(\alpha,\alpha)$ 和 $M(\alpha,\beta)$ 的值；
    \item 当 $n=4$ 时，设 $B$ 是 $A$ 的子集，且满足：对于 $B$ 中的任意元素 $\alpha,\beta$，当 $\alpha,\beta$ 相同时，$M(\alpha,\beta)$ 是奇数；当 $\alpha,\beta$ 不同时，$M(\alpha,\beta)$ 是偶数。求集合 $B$ 中元素个数的最大值；
    \item 给定不小于 $2$ 的 $n$，设 $B$ 是 $A$ 的子集，且满足：对于 $B$ 中的任意两个不同的元素 $\alpha,\beta$，$M(\alpha,\beta)=0$。写出一个集合 $B$，使其元素个数最多，并说明理由。
\end{enumerate}
    \end{enumerate}
\end{description}
\clearpage


\subsection{2018年北京卷（文）}\label{gk:2018年北京卷（文）}

\begin{description}
    \item[一、]
    \begin{enumerate}
        \item 1
        \item 2
        \item 3
        \item 4
        \item 5
        \item 6
        \item 7
        \item 8
    \end{enumerate}
    \item[二、]
    \begin{enumerate}
      \addtocounter{enumi}{+8}   
        \item 1
        \item 2
        \item 3
        \item 4
        \item 1
        \item 6
    \end{enumerate}
    \item[三、]   
    \begin{enumerate}
      \addtocounter{enumi}{+14}   
        \item 1
        \item 2
        \item 3
        \item 
        \item 设函数 $f(x)=[a x^{2}-(3 a+1) x+3 a+2]  e^{x}$。

\begin{enumerate}[label=(\arabic*)]
    \item 若曲线 $y=f(x)$ 在点 $(2,f(2))$ 处的切线斜率为 $0$，求 $a$；
    \item 若 $f(x)$ 在 $x=1$ 处取得极小值，求 $a$ 的取值范围。
\end{enumerate}
        \item 已知椭圆 $M:\frac{x^{2}}{a^{2}}+\frac{y^{2}}{b^{2}}=1(a>b>0)$ 的离心率为 $\frac{\sqrt{6}}{3}$，焦距为 $2\sqrt{2}$，斜率为 $k$ 的直线 $l$ 与椭圆 $M$ 有两个不同的交点 $A,B$。

\begin{enumerate}[label=(\arabic*)]
    \item 求椭圆 $M$ 的方程；
    \item 若 $k=1$，求 $|AB|$ 的最大值；
    \item 设 $P(-2,0)$，直线 $PA$ 与椭圆 $M$ 的另一个交点为 $C$，直线 $PB$ 与椭圆 $M$ 的另一个交点为 $D$。若 $C,D$ 和点 $Q(-\frac{7}{4},\frac{1}{4})$ 共线，求 $k$。
\end{enumerate}
    \end{enumerate}
\end{description}
\clearpage


\subsection{2018年江苏卷}\label{gk:2018年江苏卷}


\clearpage


\subsection{2018年上海卷（理）}\label{gk:2018年上海卷（理）}


\clearpage


\subsection{2018年上海卷（文）}\label{gk:2018年上海卷（文）}


\clearpage


\subsection{2018年天津卷（理）}\label{gk:2018年天津卷（理）}


\clearpage


\subsection{2018年天津卷（文）}\label{gk:2018年天津卷（文）}


\clearpage


\subsection{2018年浙江卷}\label{gk:2018年浙江卷}


\clearpage